\documentclass[11pt,a4paper]{article}
\usepackage[T1]{fontenc}
\usepackage{lmodern}
\usepackage{amsmath,amssymb}
\usepackage{booktabs}
\usepackage{array}
\usepackage[a4paper,margin=2.2cm]{geometry}
\usepackage[hidelinks]{hyperref}
\setlength{\parskip}{3pt}
\title{\vspace{-1.4cm}\Large Couplings one at a time and the closure: milestone report for Task~2}
\author{Dr.\ Vipin\\ Gurugram University}
\date{30 October 2026}
\begin{document}
\maketitle
\vspace{-7mm}

\section*{What was already settled}

The fourteen-row matrix is not singular; the earlier verdict came from scanning the determinant, a
coarse grid and the wrong length ratio. Your three dips are reproduced to the grid: A0 at
$2340.0$, $2365.0$ and $2460.3$~m/s for $l_g/h=0.05$, $0.10$, $0.20$ at $h=5\times10^{-4}$~m and
$fd=1000$, with depths of $92$, $33$ and $63$ against the local median. The gradient length
stiffens both branches monotonically: S0 runs $5311.2$, $5318.7$, $5348.4$ at the same three
lengths against $5308.6$ elastic, and A0 runs $2337.4$, $2367.0$, $2460.3$ against $2326.6$. The
$fd=10$ S0 entries stay excluded for the reason already accepted.

\section*{The piezoelectric limit, corrected}

You were right that the $5368$~m/s figure was not a fault. The carriers were still on and mobile
charge was screening the stiffening. With $q=0$, the thermal couplings off, $p_3=0$ and
$b_1=b_3=0$, the open-circuit thin-plate speed comes out where it should:

\begin{center}\footnotesize
\begin{tabular}{lccc}
\toprule
case & $fd=10$ & $fd=100$ & closed form\\
\midrule
piezoelectric, faces open & 5509.7 & 5509.2 & 5509.57\\
elastic baseline, $e=0$ & 5259.5 & 5258.8 & 5259.13\\
\bottomrule
\end{tabular}
\end{center}

That is 0.06 per cent at worst, and it confirms the plane-stress reduction the masters carry.

Switching the carriers back on gives the screening curve at $fd=100$: $5509.2$ at $n_0=0$,
$5509.2$ at $10^{16}$, $5504.9$ at $10^{18}$, $5474.7$ at $10^{19}$, $5355.1$ at $10^{20}$ and
$5273.7$~m/s at $10^{21}$~m$^{-3}$. The crossover falls between $10^{18}$ and $10^{19}$, and at the
top of the range the speed has come down to the short-circuit value, which is what complete
screening means. This curve is worth a figure in the Sumit paper.

\section*{The closure, with the two lengths separated}

Your diagnosis was exact. With $l_{nl}=l_g$ the two operators coincide, $N(\eta)=G(\eta)$, the
elastic block becomes $G$ times the classical one, and at the four gradient roots $G$ vanishes, so
the traction rows divide by zero. That is the origin of the $10^{-18}$ singular values, not the
closure. With $l_{nl}=0.5\,l_g$ and $l_{nl}=2\,l_g$ the runs are healthy: medians of $10^{-2}$ to
$10^{-1}$ and dips of depth 87 to 420.

Sumit paper fully coupled, $h=5\times10^{-4}$~m, $fd=1000$, A0 branch, elastic value $2326.6$~m/s:

\begin{center}\footnotesize
\begin{tabular}{lccc}
\toprule
$l_g/h$ & $l_{nl}/l_g$ & our seven rows & alternative closure\\
\midrule
0.05 & 0.5 & 2193.2 & 2211.5\\
0.05 & 2 & 2175.8 & 2170.8\\
0.10 & 0.5 & 2220.2 & 2195.3\\
0.10 & 2 & 2154.4 & 2128.4\\
0.20 & 0.5 & 2300.7 & no dip near 2300\\
0.20 & 2 & 1825 & 1825\\
\bottomrule
\end{tabular}
\end{center}

Where both closures give a dip in the same window the difference is at most 26~m/s, 1.1 per cent,
and carries no systematic sign. That is smaller than the size effect it is competing with, since
these lengths move the branch by 3 to 22 per cent, so the closure does not decide the numbers. It
does decide what we may claim: the seven rows remain a modelling choice of ours, and the masters
still need the paragraph that says so. The proposed wording now carries the 1.1 per cent figure.

Two entries need care. Every value in the table lies below the elastic speed, even with $l_{nl}$ at
half the gradient length, because at these lengths the nonlocal softening outweighs the gradient
stiffening; that ordering should be stated where the limits are quoted. And the
$l_{nl}=0.5\,l_g$, $l_g/h=0.2$ alternative closure has no dip anywhere near 2300~m/s but a shallow
one at 1110~m/s. A closure that removes a branch is not credible, so that entry is held back and
will be re-run; the likely cause is the same near-degenerate column pair, now appearing inside one
of the two row sets at the larger length.

\section*{Parity}

All thirty-two diagonal sign vectors were tested against $M(-\eta)=S\,M(\eta)\,S$ at
$k=3\times10^{5}$, $\omega=2\pi10^{8}$ and $\eta=10^{4}+2\times10^{3}i$.

\begin{center}\footnotesize
\begin{tabular}{llc}
\toprule
paper & field order & patterns\\
\midrule
Sumit & $u_x,u_z,\phi,\psi,\theta$ & none\\
Sumit, $\beta_{th}=\Pi_{th}=E_g=0$ & same & $(+,-,-,-,+)$ and its negative\\
Neha & $u_x,u_z,w,\phi,\theta$ & $(+,-,-,-,+)$ and its negative\\
Anu & $u_x,u_z,\phi,\psi,\theta$ & $(+,-,-,-,+)$ and its negative\\
\bottomrule
\end{tabular}
\end{center}

This is the wording the three masters already carry, and it stands as written.

\section*{Fourier damping: attempted, not a test yet}

A bulk longitudinal wave at $\eta=0$ with the thermal coupling live and everything else off gives
$k_i/k_r$ of order $10^{-19}$ at $100$~kHz and $10^{-14}$ at $10$~GHz, rising straight through the
frequency where the Fourier shape should peak, $1/(2\pi\tau_R)=159$~MHz. Attenuation five to seven
orders below any physical $\Delta_E$ cannot be right, and the likely reason is that the defaults
mingle units: $bth=10^{18}$ cannot be SI beside $b_1=b_3=4\times10^{5}$. The damping check, and with
it the Hutson--White peak, the phason limits and the magnetoelastic limits, waits until the two are
put on one footing. I would rather say that than quote a number.

\section*{Delivery}

The v8 masters are untouched. Every number here is in the internal notes with the code that
produced it. No \LaTeX{} engine is installed in this sandbox, so this file is a source; it has been
through the style script and compiles on any recent \TeX{}~Live.

\end{document}
